\chapter{Trechos do código-fonte}\label{ap:codigo}

Os trechos abaixo foram copiados sem alteração da versão 3D (pasta \arquivo{Sim3D\_2.0/src}). Os comentários originais foram mantidos porque fazem parte da documentação do projeto: muitos explicam \emph{por que} uma decisão foi tomada, e não apenas o que o código faz. Os arquivos completos estão na pasta \arquivo{codigo/} que acompanha este relatório.

\lstinputlisting[language=TypeScript,firstline=34,lastline=71,
  caption={Aceleração gravitacional mútua (\arquivo{nucleo/forcas.ts})},label=cod:forcas]{codigo/forcas.ts}

\lstinputlisting[language=TypeScript,firstline=65,lastline=91,
  caption={Integrador de Yoshida de quarta ordem (\arquivo{nucleo/integrador.ts})},label=cod:yoshida]{codigo/integrador.ts}

\lstinputlisting[language=TypeScript,firstline=56,lastline=103,
  caption={Solução da equação de Kepler (\arquivo{nucleo/kepler.ts})},label=cod:kepler]{codigo/kepler.ts}

\lstinputlisting[language=TypeScript,firstline=133,lastline=178,
  caption={Elementos orbitais para posição e velocidade (\arquivo{nucleo/kepler.ts})},label=cod:elementos]{codigo/kepler.ts}

\lstinputlisting[language=TypeScript,firstline=91,lastline=131,
  caption={Coeficientes de Fourier complexos (\arquivo{nucleo/fourier.ts})},label=cod:decompor]{codigo/fourier.ts}

\lstinputlisting[language=TypeScript,firstline=182,lastline=221,
  caption={Cadeia de epiciclos (\arquivo{nucleo/fourier.ts})},label=cod:cadeia]{codigo/fourier.ts}

\lstinputlisting[language=TypeScript,firstline=1,lastline=38,
  caption={Origem flutuante (\arquivo{cena/origemFlutuante.ts})},label=cod:origem]{codigo/origemFlutuante.ts}

\lstinputlisting[language=TypeScript,firstline=26,lastline=55,
  caption={Interpolação de Hermite para o desenho (\arquivo{nucleo/amostragem.ts})},label=cod:hermite]{codigo/amostragem.ts}
